\subsubsection{Split Dataset}
    A customized under-sampling technique was used to avoid this situation. 
    The standard under-sampling technique does not meet the requirement simply due to the fact that many of the images used for the dataset have multiple labels.
    In other words, elimination of an image will lead to the perishment of many other labels.
    \paragraph{Calculating Quantity}
        Firstly, start with finding the lowest number of labels out of all the classes to set a standard upon which to divide a ratio between training, validating, and testing. The sum of the ratio of each dataset must be equal to 1:\\
        \begin{center}
            \(p(train) + p(val) + p(test) = 1\) 
        \end{center} 
        From that, calculate the number of labels per class needed in each subset by:
        \begin{center}
            \(n(train) = p(train) * standard\) \\
            \(n(val) = p(val) * standard\) \\
            \(n(test) = p(test) * standard\) 
        \end{center} 
    \paragraph{Customized Under-sampling Algorithm}
        The choose approach was employed as opposed to deleting. 
        The algorithm will first loop over all of the classes that are available. 
        After that, sort the dataset by how many labels there are in total for that class for each image. 
        The image will be chosen by an algorithm from left to right. 
        The image will be added to the new dataset if the total number of labels from each class does not exceed the predetermined number in the previous section (\autoref{fig:CustomizedUnderSampling}).